\section{Conclusion}
\label{sec:conclusion}

This paper asked whether reclamation can be a property of the emitted program rather than a service
attached to it. The answer is yes, and in a form a user can sign: the compiler decides, per unit of
compilation, whether the unit needs a reclamation run-time at all, and when it does, the run-time it
emits satisfies two bounds that the paper states as theorems and the build checks as gates.

\paragraph{What the machine is.}
Five storage substrates --- a frame arena, a slab-allocated region with occupancy and mark bitmaps, a
container pool, per-task arenas, and an \code{mmap}-backed plate --- are enough to cover the language
and small enough to reason about.  Above them sit four release layers that collapse, per unit, into
three tiers: no run-time at all, an event tier that counts and never sweeps, and a sweep tier that
reclaims without counting.  The tier is chosen from a static acyclicity judgement, and a separate
per-function analysis removes the run-time from functions that store no escaping reference.  The
census is where the design pays for itself: 68.5\% of the programs in the corpus that compile emit no
reclamation code whatsoever, so the mechanism that carries the corpus is the analysis, not the collector.

\paragraph{What is proved.}
Seven theorems.  Safety: no execution reaches a freed slot, with the static root table as the only
structural assumption, and with the four release layers --- block exit, static placement, sweep, and
release at the store that orphans a structure --- each discharged as a lemma.  Tier soundness: the
event tier cannot collect a reachable object, and the tier switch is semantics-preserving.  Idle: a
unit judged idle can have its run-time dropped without changing the observable trace.  Memory: a
loop with a constant live set runs in \mth{L + G + O(B)}, flat in the number of iterations, which
\tab\ref{tab:bounded} shows as 3,420\,kB at 125 thousand and at eight million rounds.  Pause: a
single step is bounded by the batch, the scan chunk, the segment deadline, the sweep credit and a
constant, all of which are compile-time numbers.  Amortised work: linear in allocated bytes per
region for the acyclic tier, and linear in capacity plus live set for the sweep tier.  Progress: a
deferred free queue drains under its budget or the allocation that could not find room aborts instead
of piling up.

\paragraph{What the roofline buys.}
The four laws turn the tuning questions into geometry.  Cost: reclamation time per allocated byte
falls as \mth{c_s F + (c_t + c_s F)/(r-1)}, so the trigger ratio is the only knob that matters and
its floor is set by the slab walk.  Memory: peak heap is \mth{L + \kappa_M H}, which fixes the
largest achievable \mth{r}.  Pause: a step costs \mth{b \tau_t + W_s 64 \tau_s}, which fixes the
largest usable batch.  Keep-up: the collector must retire garbage as fast as the program makes it,
which needs \mth{r \ge 1 + gK/b}.  Feasibility is then one inequality between the memory budget and
the pause budget, and the operating point is its corner.  With six constants calibrated to two to
ten percent, the model predicted the measured cost of the \code{digraph} workload to within 1.7\%
and located the default policy at 33 times below the allocation platform: the runtime finishes its
work far inside the budget the roofline says it has.

\paragraph{What the numbers say.}
Against the best of C++ and Rust on the same four shapes, extC is 2.81 times and 1.00 times faster on
\code{promote}, 1.12 and 1.29 on \code{sizeratio}, 1.87 times faster than the best Rust arm on
\code{chase} while level with the best C++ arm there, and 2.40 and 1.19 on \code{cycleratio}; its
resident set is never larger than the best arm of any comparison language and is 13\% smaller than C++
on the shape where memory is the point.  No shape has a single step above one millisecond, the 99.9th
percentile of a step is between 10 and 59 \us, and the total time in reclamation ranges from 3.8\% of
wall clock on \code{chase} to 40.5\% on \code{promote} --- a share that a concurrent collector hides
on another thread, and that this design charges to the mutator precisely so that it can be bounded.

\paragraph{What the contracts unlock.}
Three levers, in the order we would pull them.  First, the acyclicity proof: the counting tier is
proved, sound, and removes the sweep entirely, and it fires on two of 794 corpus programs because its
freshness conditions are what a heap built by a loop rarely satisfies; widening those conditions with a
proof is the highest-leverage change on the table, because it deletes work rather than dividing it.
Second, the store barrier is the remaining cost on pointer-chasing code, and the phase analysis of
\S\ref{sec:eval} locates it in a layout choice rather than in the collector; the design space of
moving collectors that preserve the deterministic layout the language promises is open, and entering it
is the natural next question.  Third, the contract is single-threaded: a per-worker instance of the
same two bounds is what the \code{threads} block adds, and the design already gives it per worker,
since threads do not share managed regions.
